% ============================================================
%  Section 8 -- SQL: Structured Query Language
% ============================================================

\section{Structured Query Language}

% ----------------------------------------------------------
\begin{frame}{What is SQL?}
  \textbf{SQL} stands for \emph{Structured Query Language}. Originally called
  \emph{SEQUEL}, it was developed in \textbf{1974 at IBM Research} by Donald
  Chamberlin and Raymond Boyce for the System~R relational prototype.

  \begin{block}{Three Sub-Languages}
    \begin{tabular}{@{}lll@{}}
      \toprule
      \textbf{Sub-language} & \textbf{Statements} & \textbf{Purpose} \\
      \midrule
      \textbf{DML} -- Data Manipulation & \texttt{SELECT, INSERT, UPDATE, DELETE} & Work with data \\
      \textbf{DDL} -- Data Definition   & \texttt{CREATE, ALTER, DROP}            & Define schema \\
      \textbf{DCL} -- Data Control      & \texttt{GRANT, REVOKE}                  & Manage access \\
      \bottomrule
    \end{tabular}
  \end{block}

  \begin{alertblock}{Case Sensitivity}
    SQL \textbf{keywords} (\texttt{SELECT}, \texttt{FROM}, \ldots) are NOT case-sensitive.
    String \emph{values} ARE case-sensitive: \texttt{'Hamburg'} $\neq$ \texttt{'hamburg'}.
  \end{alertblock}
\end{frame}

% ----------------------------------------------------------
\begin{frame}{SQL Standardisation History}
  \begin{block}{Key Standards}
    \begin{tabular}{@{}ll@{}}
      \toprule
      \textbf{Standard} & \textbf{Key additions} \\
      \midrule
      \textbf{SQL-86}   & First ANSI/ISO standard \\
      \textbf{SQL-92}   & Major expansion; baseline for most textbooks \\
      \textbf{SQL:1999} & Recursion, triggers, object extensions \\
      \textbf{SQL:2003} & Window functions, XML support \\
      \textbf{SQL:2016} & JSON support, row pattern matching \\
      \bottomrule
    \end{tabular}
  \end{block}

  \begin{block}{Dialects}
    Every RDBMS implements the core standard but adds its own extensions:
    \begin{itemize}
      \item \textbf{MySQL/MariaDB}: \texttt{LIMIT}, \texttt{AUTO\_INCREMENT}
      \item \textbf{PostgreSQL}: rich window functions, arrays, JSON operators
      \item \textbf{Oracle}: \texttt{ROWNUM}, \texttt{CONNECT BY} (hierarchies)
      \item \textbf{SQL Server}: \texttt{TOP}, T-SQL procedural extensions
    \end{itemize}
    This course uses \textbf{MySQL 8.x} syntax throughout.
  \end{block}
\end{frame}

% ----------------------------------------------------------
\begin{frame}[fragile]{The Example Database -- Music Schema}
  All SQL examples in this section use the following schema:

  \begin{block}{Tables}
\begin{lstlisting}[style=sqlstyle]
CREATE TABLE Artist (
    ArtistID   INT PRIMARY KEY AUTO_INCREMENT,
    Name       VARCHAR(100) NOT NULL,
    Country    VARCHAR(50)
);

CREATE TABLE Album (
    AlbumID     INT PRIMARY KEY AUTO_INCREMENT,
    Title       VARCHAR(150) NOT NULL,
    ArtistID    INT NOT NULL,
    Genre       VARCHAR(50),
    ReleaseYear INT,
    FOREIGN KEY (ArtistID) REFERENCES Artist(ArtistID)
);

CREATE TABLE Track (
    TrackID   INT PRIMARY KEY AUTO_INCREMENT,
    Title     VARCHAR(150) NOT NULL,
    AlbumID   INT NOT NULL,
    Duration  INT,           -- duration in seconds
    FOREIGN KEY (AlbumID) REFERENCES Album(AlbumID)
);
\end{lstlisting}
  \end{block}
\end{frame}

% ----------------------------------------------------------
\begin{frame}[fragile]{The Example Database -- Sample Data}
  \begin{block}{Sample Rows}
\begin{lstlisting}[style=sqlstyle]
-- Artists
INSERT INTO Artist VALUES (1, 'The Beatles',  'UK');
INSERT INTO Artist VALUES (2, 'Kraftwerk',    'Germany');
INSERT INTO Artist VALUES (3, 'Daft Punk',    'France');

-- Albums
INSERT INTO Album VALUES (1, 'Abbey Road',      1, 'Rock',        1969);
INSERT INTO Album VALUES (2, 'Autobahn',        2, 'Electronic',  1974);
INSERT INTO Album VALUES (3, 'Random Access Memories', 3, 'Electronic', 2013);

-- Tracks (Duration in seconds)
INSERT INTO Track VALUES (1, 'Come Together', 1, 259);
INSERT INTO Track VALUES (2, 'Something',     1, 182);
INSERT INTO Track VALUES (3, 'Autobahn',      2, 1373);
INSERT INTO Track VALUES (4, 'Get Lucky',     3, 248);
INSERT INTO Track VALUES (5, 'Instant Crush', 3, 337);
\end{lstlisting}
  \end{block}
\end{frame}
